\newpage
\phantomsection
\addcontentsline{toc}{section}{附录}
\section*{附~~录}
\vspace*{-2mm}
\zihao{5}
\setlength{\baselineskip}{25pt}

本附录给出系统核心模块的完整源代码，包括预处理流水线、FBCSP 特征提取器和在线仿真服务端。所有代码均采用 Python 3.10+ 编写，依赖库版本见表\ref{tab:tech-stack}。

% =====================================================================
\subsection*{A\quad 预处理流水线}
% =====================================================================

代码\ref{code:pipeline}为预处理流水线的主控模块，按照第四章所述的 7 步流程依次执行坏导检测、共平均参考、带通滤波、陷波滤波、ICA 伪迹去除、事件分段和坏段拒绝。各步骤的开关与参数均由 YAML 配置文件驱动（见第\ref{sec:config-system}节）。

\begin{lstlisting}[language=Python, caption={预处理流水线 \texttt{src/preprocessing/pipeline.py}\label{code:pipeline}}]
class PreprocessingPipeline:
    """Config-driven preprocessing pipeline."""

    def __init__(self, config: dict[str, Any]):
        self.config = config
        self._ica: mne.preprocessing.ICA | None = None
        self._quality_report: dict | None = None

    def run(self, raw: BaseRaw, events: np.ndarray,
            event_id: dict[str, int]) -> mne.Epochs:
        cfg = self.config

        # 1. Bad channel detection & interpolation
        if cfg.get("bad_channel_detection", False):
            detect_bad_channels(raw)
            interpolate_bad_channels(raw)

        # 2. Common Average Reference
        ref = cfg.get("reference")
        if ref == "car":
            set_common_average_reference(raw)

        # 3. Bandpass filter
        bp = cfg.get("bandpass", [8, 30])
        bandpass_filter(raw, l_freq=bp[0], h_freq=bp[1])

        # 4. Notch filter (optional)
        notch = cfg.get("notch_filter")
        if notch:
            notch_filter(raw, freqs=notch)

        # 5. ICA artifact removal (optional)
        art_cfg = cfg.get("artifact_rejection", {})
        if art_cfg.get("method") == "ica":
            raw, self._ica = run_ica_artifact_removal(
                raw,
                n_components=art_cfg.get("n_components", 15),
                eog_channels=art_cfg.get("eog_channels"),
            )

        # 6. Epoching
        ep_cfg = cfg.get("epoch", {})
        tmin = ep_cfg.get("tmin", -0.5)
        tmax = ep_cfg.get("tmax", 3.0)
        baseline = ep_cfg.get("baseline")
        if isinstance(baseline, list):
            baseline = tuple(baseline)
        epochs = create_epochs(raw, events, event_id,
                               tmin=tmin, tmax=tmax,
                               baseline=baseline)

        # 7. Bad epoch rejection
        threshold = ep_cfg.get("reject_threshold")
        if threshold is not None:
            epochs = reject_bad_epochs(epochs,
                                       reject_threshold=threshold)

        self._quality_report = generate_quality_report(
            raw, epochs, self._ica)
        return epochs
\end{lstlisting}

% =====================================================================
\subsection*{B\quad FBCSP 特征提取器}
% =====================================================================

代码\ref{code:fbcsp}为 FBCSP 特征提取器的完整实现，对应第五章第\ref{sec:fbcsp-theory}节的算法描述。\texttt{fit()} 方法执行子频带分解、逐带 CSP 和互信息特征选择三个步骤；\texttt{transform()} 方法将新数据通过已拟合的滤波器组进行特征投影。

\begin{lstlisting}[language=Python, caption={FBCSP 特征提取器 \texttt{src/features/fbcsp.py}\label{code:fbcsp}}]
class FBCSPExtractor:
    """Filter Bank CSP feature extractor."""

    def __init__(self, filter_bands=None, n_components=4,
                 selection_method="mutual_info",
                 n_select=8, sfreq=250.0):
        self.filter_bands = filter_bands or [
            (4, 8), (8, 12), (12, 16), (16, 20),
            (20, 24), (24, 28), (28, 32),
        ]
        self.n_components = n_components
        self.n_select = n_select
        self.sfreq = sfreq
        self._csp_list: list[CSP] = []
        self._selected_indices: list[int] | None = None

    def fit(self, X: np.ndarray, y: np.ndarray):
        self._csp_list = []
        all_features = []
        n_channels = X.shape[1]
        actual_comp = min(self.n_components, n_channels - 1)
        if actual_comp < 2:
            actual_comp = 2
        self._actual_components = actual_comp
        reg = "ledoit_wolf" if n_channels <= 4 else None

        for lo, hi in self.filter_bands:
            X_filt = filter_data(X.astype(np.float64),
                                 sfreq=self.sfreq,
                                 l_freq=lo, h_freq=hi,
                                 verbose=False)
            csp = CSP(n_components=actual_comp, reg=reg,
                      log=True, norm_trace=False)
            feats = csp.fit_transform(X_filt, y)
            self._csp_list.append(csp)
            all_features.append(feats)

        X_concat = np.hstack(all_features)
        actual_select = min(self.n_select, X_concat.shape[1])
        _, self._selected_indices = mutual_information_selection(
            X_concat, y, n_select=actual_select)
        return self

    def transform(self, X: np.ndarray) -> np.ndarray:
        all_features = []
        for (lo, hi), csp in zip(self.filter_bands,
                                  self._csp_list):
            X_filt = filter_data(X.astype(np.float64),
                                 sfreq=self.sfreq,
                                 l_freq=lo, h_freq=hi,
                                 verbose=False)
            all_features.append(csp.transform(X_filt))
        X_concat = np.hstack(all_features)
        return X_concat[:, self._selected_indices]

    def get_selected_band_info(self) -> list[dict]:
        nc = getattr(self, "_actual_components",
                     self.n_components)
        info = []
        for idx in (self._selected_indices or []):
            band_idx = idx // nc
            info.append({
                "feature_index": idx,
                "band_range": self.filter_bands[band_idx],
                "csp_component": idx % nc,
            })
        return info
\end{lstlisting}

% =====================================================================
\subsection*{C\quad 在线仿真 WebSocket 服务端}
% =====================================================================

代码\ref{code:server}为在线仿真系统的 WebSocket 服务端核心逻辑，对应第七章的系统架构。\texttt{\_run\_trials()} 方法实现了"提示 → 想象 → 分类 → 反馈 → 休息"的试次状态机，通过 \texttt{\_broadcast()} 将 JSON 状态消息推送给所有已连接的浏览器客户端。

\begin{lstlisting}[language=Python, caption={在线仿真服务端 \texttt{src/online/server.py}\label{code:server}}]
TIMING = {"cue": 2.0, "imagine": 3.0,
          "feedback": 2.0, "rest": 1.5}

class BCIServer:
    """WebSocket server driving a BCI trial loop."""

    def __init__(self, classifier, replay, host="localhost",
                 port=8765, http_port=8080, max_trials=None):
        self.classifier = classifier
        self.replay = replay
        self.host, self.port = host, port
        self.http_port = http_port
        self.max_trials = max_trials
        self._clients: set = set()
        self._running = False

    async def _broadcast(self, data: dict) -> None:
        msg = json.dumps(data, ensure_ascii=False)
        dead = set()
        for ws in self._clients:
            try:
                await ws.send(msg)
            except websockets.ConnectionClosed:
                dead.add(ws)
        self._clients -= dead

    async def _run_trials(self) -> None:
        self._running = True
        trial_num, correct = 0, 0

        for epoch, true_label in self.replay:
            if not self._running:
                break
            if self.max_trials and trial_num >= self.max_trials:
                break
            trial_num += 1
            true_name = LABEL_MAP.get(true_label,
                                       str(true_label))
            # CUE
            await self._broadcast({
                "state": "cue", "trial_num": trial_num,
                "cue_direction": true_name})
            await asyncio.sleep(TIMING["cue"])

            # IMAGINE
            await self._broadcast({
                "state": "imagine", "trial_num": trial_num,
                "cue_direction": true_name})
            await asyncio.sleep(TIMING["imagine"])

            # CLASSIFY
            result = self.classifier.predict_trial(epoch)
            is_correct = result["label"] == true_label
            if is_correct:
                correct += 1

            # FEEDBACK
            await self._broadcast({
                "state": "feedback",
                "trial_num": trial_num,
                "prediction": result["label"],
                "confidence": result["confidence"],
                "correct": is_correct,
                "latency_ms": result["latency_ms"],
                "accuracy": round(correct / trial_num, 4)})
            await asyncio.sleep(TIMING["feedback"])

            # REST
            await self._broadcast({
                "state": "rest", "trial_num": trial_num})
            await asyncio.sleep(TIMING["rest"])

        # Session complete
        await self._broadcast({
            "state": "done", "total_trials": trial_num,
            "total_correct": correct,
            "final_accuracy": round(
                correct / max(trial_num, 1), 4)})

    async def run(self) -> None:
        self._start_http_server()
        await websockets.serve(self._ws_handler,
                               self.host, self.port)
        await asyncio.Future()  # run forever
\end{lstlisting}

% =====================================================================
\subsection*{D\quad 默认实验配置文件}
% =====================================================================

代码\ref{code:config}为系统的全局默认 YAML 配置文件，定义了预处理、特征提取、分类模型和评估策略的所有可配置参数。实验配置通过三层合并机制（默认 → 数据集 → 实验）覆盖相应字段，详见第\ref{sec:config-system}节。

\begin{lstlisting}[caption={全局默认配置 \texttt{configs/default.yaml}\label{code:config}}]
seed: 42

preprocessing:
  reference: "car"
  bandpass: [8, 30]
  notch_filter: null
  bad_channel_detection: false
  artifact_rejection:
    method: null
    n_components: 15
    eog_channels: null
  epoch:
    tmin: -0.5
    tmax: 3.0
    train_tmin: 0.5
    train_tmax: 2.5
    baseline: null
    reject_threshold: 100.0e-6

features:
  method: "csp"
  params:
    n_components: 4
    log: true

model:
  type: "lda"
  params: {}

evaluation:
  method: "stratified_kfold"
  n_splits: 10
  metrics: ["accuracy", "kappa", "f1_weighted"]

output:
  save_model: true
  save_figures: true
  save_replay_data: true
\end{lstlisting}


\label{appendix}  % 请勿删除此行
